\chapter{Hardware Deployment and Preliminary Experimental Results}
\label{ch:hardware}

% ---------------------------------------------------------------------
\section{Deployment on the \hero{} Platform}
\label{sec:hw-deployment}

% TODO: startup sequence (Zenoh routers paired -> host MPC -> IPC MRT
% bridge); wired vs. wireless link (wired chosen; must be stated as a
% run condition); first successful real-robot run milestone.

% ---------------------------------------------------------------------
\section{Deployment Challenges and Mitigations}
\label{sec:hw-challenges}

% TODO:
%  - policy lag: solve time ~147-225 ms (mean), effective MPC rate
%    ~4.2 Hz vs mpcDesiredFrequency=10; policyLag ~0.37 s (down 2.2x
%    from ~0.8 s); DDP Initialization ~55.5% of solve time.
%  - solve-time spikes (>1000 ms) vs ~1.03 s policy window -> bridge
%    enters hold; mitigation = wider solutionTimeWindow (1.0 s),
%    periodic re-publish, diagnostics ([policy_interp]/[policy_arrival]).
%  - re-handshake silently overwriting reference -> 1 Hz re-publish.
%  - cabling tethering the base (odom contamination) after switching to
%    wired link.

% ---------------------------------------------------------------------
\section{Preliminary End-Effector Tracking Results}
\label{sec:hw-results}

% TODO: EE-pose tracking on the four-point rectangle; RMS/max per
% straight segment and per corner; tangential lag ~ policyLag x v_ref
% (H2, ~3.7 cm expected); closed-loop drift; per-sample join of
% policyLag vs lag_eff. State the odom+FK internal-metrology limitation.
